\section{Минимизация булевых функций. Минимальные и сокращенные ДНФ. Геометрическая интерпретация.}

\subsection*{Постановка задачи минимизации}

Булева функция может быть реализована бесконечным множеством равносильных формул. Задача минимизации состоит в поиске самой «простой» формулы.
\paragraph{Критерии простоты:}
\begin{itemize}
    \item Наименьшее количество вхождений переменных.
    \item Наименьшее количество логических операций.
\end{itemize}
Наиболее изучена задача минимизации в классе **дизъюнктивных нормальных форм (ДНФ)**.

\subsection*{Дизъюнктивные формы}

\paragraph{Определения}
\begin{itemize}
    \item \textbf{Элементарная конъюнкция} $K$ — это конъюнкция переменных или их отрицаний, в которую каждая переменная входит не более одного раза.
    \item \textbf{Ранг конъюнкции} $|K|$ — количество переменных, входящих в неё.
    \item \textbf{ДНФ} — это дизъюнкция элементарных конъюнкций: $\bigvee_{i=1}^k K_i$.
\end{itemize}
Число различных ДНФ от $n$ переменных равно $2^{3^n}$, что делает полный перебор для минимизации практически невозможным уже при $n > 5$.

\subsection*{Геометрическая интерпретация}

Новиков предлагает рассматривать наборы значений переменных как вершины **$n$-мерного единичного гиперкуба** $E^n$.
\begin{itemize}
    \item \textbf{Вершина} (0-мерная грань) соответствует набору из $n$ констант (элементарной конъюнкции ранга $n$).
    \item \textbf{Ребро} (1-мерная грань) соответствует конъюнкции ранга $n-1$ (две вершины, различающиеся в одном разряде).
    \item \textbf{$k$-мерная грань} соответствует конъюнкции ранга $n-k$. Чем меньше переменных в конъюнкции, тем большую область («облако» единиц) в кубе она покрывает.
\end{itemize}
Булева функция задается подмножеством вершин гиперкуба, на которых она равна $1$. Задача минимизации превращается в задачу **покрытия всех «единичных» вершин минимальным числом граней максимальной размерности**.

\subsection*{Сокращенные и минимальные ДНФ}

\paragraph{Простая импликанта}
Элементарная конъюнкция $K$ называется **импликантой** функции $f$, если $K=1 \implies f=1$. Импликанта называется **простой**, если при удалении из неё любой переменной она перестает быть импликантой (геометрически: это грань максимальной размерности, целиком состоящая из единиц функции).

\paragraph{Сокращенная ДНФ}
ДНФ, состоящая из дизъюнкции **всех** простых импликант функции, называется **сокращенной**.
\textit{Способ получения:}
1. Взять СДНФ функции.
2. Применять закон \textbf{склеивания}: $(x \land \mathcal{A}) \lor (\neg x \land \mathcal{A}) = \mathcal{A}$.
3. Применять закон \textbf{поглощения}: $\mathcal{A} \lor (\mathcal{A} \land \mathcal{B}) = \mathcal{A}$.

\paragraph{Минимальная ДНФ (МДНФ)}
Это ДНФ, содержащая минимальное суммарное число вхождений переменных. МДНФ всегда строится из некоторого подмножества простых импликант, входящих в сокращенную ДНФ.
\textit{Алгоритм (вкратце):} Из сокращенной ДНФ нужно выбрать такой минимальный набор конъюнкций, который в сумме «покрывает» все те же единицы, что и исходная функция.

\paragraph{Пример}
Для функции $f(x, y) = (x \land y) \lor (x \land \neg y) \lor (\neg x \land y)$:
1. Склеиваем $(x \land y)$ и $(x \land \neg y) \to x$.
2. Склеиваем $(x \land y)$ и $(\neg x \land y) \to y$.
3. Сокращенная ДНФ: $x \lor y$. Она же в данном случае будет минимальной.